%! AP Statistics — Unit 1, Lesson 1.6 — Homework
\documentclass[10pt]{article}
\usepackage{apstats-article}
\usepackage{apstats-boxes}

%=============================================================================
\begin{document}

\pageheader{Unit 1, Lesson 1.6}{Homework: Describing the Distribution of a Quantitative Variable}
\namedateperiod

\vspace{0.1in}

\begin{tcolorbox}[colback=sky, colframe=navy, boxrule=0.8pt, arc=2mm,
    left=3mm, right=3mm, top=1.5mm, bottom=1.5mm]
\small For all written descriptions: use the SOCS framework; include context
(variable name and population) in every sentence; use complete sentences.
AP free-response graders look for all four components and context throughout.
\end{tcolorbox}

\vspace{0.14in}

% ════════════════════════════════════════════════════════════════════════════
% PART A
% ════════════════════════════════════════════════════════════════════════════
\begin{blurbbox}[Part A \quad Identify the Shape]{navylight}
\small

For each distribution described below, identify the most likely shape from
this list: \textit{symmetric, skewed right, skewed left, bimodal, uniform}.
Give one piece of evidence from the description that supports your answer.

\vspace{10pt}

\begin{enumerate}[leftmargin=1.5em, itemsep=14pt]

  \item The prices of homes sold in a county last year. Most homes sold for
        \$180{,}000--\$350{,}000, but a small number of luxury estates sold
        for over \$2 million.\\[4pt]
        Shape: \blank{3.5cm} \quad Evidence: \blank{6cm}\\[1pt]
        \writeline

  \item Scores on a very easy quiz (out of 10) for a class of 30 students.
        The distribution shows most students scored 9 or 10, with only a few
        scoring below 7.\\[4pt]
        Shape: \blank{3.5cm} \quad Evidence: \blank{6cm}\\[1pt]
        \writeline

  \item The number of hours of sleep per night for 60 adults, with a peak
        near 7 hours and roughly equal drop-off on both sides.\\[4pt]
        Shape: \blank{3.5cm} \quad Evidence: \blank{6cm}\\[1pt]
        \writeline

  \item Survey responses on a 5-point scale where all responses (1, 2, 3, 4,
        and 5) were chosen by nearly equal numbers of respondents.\\[4pt]
        Shape: \blank{3.5cm} \quad Evidence: \blank{6cm}\\[1pt]
        \writeline

\end{enumerate}

\end{blurbbox}

\vspace{0.14in}

% ════════════════════════════════════════════════════════════════════════════
% PART B
% ════════════════════════════════════════════════════════════════════════════
\begin{blurbbox}[Part B \quad Write a Complete SOCS Description]{greenacc}
\small

A researcher records the \textbf{number of minutes per day} that 45 middle
school students spend on physical activity. Key features from the histogram:

\begin{itemize}[itemsep=3pt]
  \item Shape: skewed right; most students are active 20--60 minutes per day.
  \item A long tail extends to the right: a few students report 90--120 minutes.
  \item One student reported 150 minutes --- far above the rest of the group.
  \item The center of the distribution is approximately 38 minutes.
  \item Values range from 5 to 150 minutes.
\end{itemize}

\vspace{10pt}

\begin{enumerate}[leftmargin=1.5em, itemsep=10pt]

  \item Write a complete SOCS description of the distribution of daily physical
        activity time for these 45 middle school students. Use complete sentences
        with context in every sentence. Label each sentence S, O, C, or S in
        the margin.\\[6pt]
        \writeline\\[1pt]\writeline\\[1pt]\writeline\\[1pt]\writeline\\[1pt]
        \writeline\\[1pt]\writeline

  \item If the one student reporting 150 minutes is removed, predict how (if
        at all) the shape, center, and spread of the distribution would change.
        Explain your reasoning.\\[6pt]
        \writeline\\[1pt]\writeline\\[1pt]\writeline

  \item Based on the shape of this distribution, would you expect the mean
        daily activity time to be greater than, equal to, or less than the
        median? Explain why.\\[6pt]
        \writeline\\[1pt]\writeline\\[1pt]\writeline

\end{enumerate}

\end{blurbbox}

\vspace{0.14in}

% ════════════════════════════════════════════════════════════════════════════
% PART C
% ════════════════════════════════════════════════════════════════════════════
\begin{blurbbox}[Part C \quad AP-Style Free Response]{redacc}
\small

The following summary describes the distribution of \textbf{fuel efficiency
(miles per gallon)} for a random sample of 60 cars tested by a consumer group.

\vspace{8pt}
\renewcommand{\arraystretch}{1.5}
\begin{tabularx}{0.72\linewidth}{@{} >{\bfseries}p{0.30\linewidth}
                                      X @{}}
\rowcolor{sky}
\textbf{Feature} & \textbf{Value / Description} \\
\midrule
Shape & Roughly symmetric, unimodal \\
Minimum & 18 mpg \\
Maximum & 47 mpg \\
Approximate center & 32 mpg \\
Outliers & One car at 47 mpg; no low-end outliers \\
\end{tabularx}

\vspace{10pt}

\begin{enumerate}[leftmargin=1.5em, itemsep=10pt]

  \item Describe the distribution of fuel efficiency for these 60 cars using
        the SOCS framework. Write in complete sentences and include context in
        every sentence.\\[6pt]
        \writeline\\[1pt]\writeline\\[1pt]\writeline\\[1pt]\writeline

  \item The consumer group adds 5 more cars to the sample; all 5 have fuel
        efficiency of 14 mpg. Predict how this would change each SOCS component.\\[6pt]
        \writeline\\[1pt]\writeline\\[1pt]\writeline

  \item A classmate says: ``Since this distribution is roughly symmetric,
        the mean and median fuel efficiency should be approximately equal.''
        Do you agree? Explain.\\[6pt]
        \writeline\\[1pt]\writeline

\end{enumerate}

\end{blurbbox}

\vspace{0.14in}

% ════════════════════════════════════════════════════════════════════════════
% EXTENSION
% ════════════════════════════════════════════════════════════════════════════
\begin{extensionbox}
\small
\textbf{Optional Extension --- Real Data SOCS Description}\\[4pt]
Find a real histogram, dotplot, or stemplot in a news article, textbook, or
data website (e.g., census.gov, gapminder.org). Write a complete SOCS
description of the distribution. Then explain: what does the shape of this
distribution suggest about the real-world process that generated it?

\vspace{8pt}
\writeline\\[1pt]\writeline\\[1pt]\writeline\\[1pt]\writeline\\[1pt]\writeline

\vspace{4pt}
\textcolor{slate}{\scriptsize Source / URL:\;\blank{11cm}}
\end{extensionbox}

\end{document}
